\section{Theory Details for Proposition~\ref{prop:endpoint-energy-error}}
\label{app:theory}
\label{app:one-reset-error}

\newtheorem{corollary}{Corollary}[section]

\subsection{From Factored Dynamics to Local Affine Representation Dynamics}
\label{subsec:local-affine-dynamics}

We first justify the switched affine representation model used in the main paper.
The key point is not that manipulation dynamics are globally linear.
Rather, the claim is that after factorization, the local dependency structure of the target factor changes across passive and interaction modes.
This change forces us to distinguish the passive operator $A_0$ from the interaction operator $A_1$, and to include an interaction-dependent affine term $b_t$.

Let the state be factored as $s_t=(u_t,y_t)$, where $u_t$ is the actuated factor and $y_t$ is the target factor.
We write the target representation as
\[
    \psi_t \triangleq \psi(y_t).
\]
The CRL critic conditions on the anchor through $\phi(s_t,a_t)=\phi(u_t,y_t,a_t)$.
If the action is produced by a deterministic policy, $a_t=\pi(u_t,y_t,g)$, then the same argument can be viewed as conditioning on $(u_t,y_t)$ under the current policy.
Thus, for the theory, we only need to study how the anchor variables induce the next target representation $\psi_{t+1}$.

\paragraph{Dependency contrast.}
In locomotion-like domains, the controlled factor and the goal factor are typically the same state variable.
If this factor is denoted by $z_t$, the local dynamics have the form
\[
    z_{t+1}=f_{\rm loc}(z_t,a_t),
\]
and the action affects the goal-relevant state almost everywhere:
\[
    \left\|
    \frac{\partial f_{\rm loc}(z_t,a_t)}{\partial a_t}
    \right\|
    >0
    \quad
    \text{for a.e. } (z_t,a_t).
\]
Under a fixed policy or fixed anchor action, this gives a single smooth closed-loop transition map.
This is the setting in which a single local temporal operator is a natural approximation for the future representation, consistent with the linear-Gaussian view of temporal contrastive representations~\citep{eysenbach2024inference}.

Manipulation has a different dependency structure.
The target factor is often not directly actuated.
Using the notation from the main paper, the target dynamics are
\[
y_{t+1}
=
\begin{cases}
f_0(y_t), & \omega_t=0,\\
f_1(u_t,y_t,a_t), & \omega_t=1,
\end{cases}
\]
where $\omega_t=0$ denotes passive target dynamics and $\omega_t=1$ denotes an interaction.
Within each smooth mode, the action-to-target derivative is therefore gated by the interaction indicator:
\[
    \frac{\partial y_{t+1}}{\partial a_t}
    =
    \begin{cases}
    0, & \omega_t=0,\\
    \frac{\partial f_1(u_t,y_t,a_t)}{\partial a_t}, & \omega_t=1.
    \end{cases}
\]
Thus, action influence on the target is absent in the passive mode and appears only during interaction.
This is a structural property of manipulation, not a choice of parameterization.

\begin{assumption}[Mode-dependent target controllability]
\label{ass:mode-dependent-controllability}
Let
\[
    H_0(y) \triangleq \psi(f_0(y)),
    \qquad
    H_1(u,y,a) \triangleq \psi(f_1(u,y,a)).
\]
We assume that $H_0$ and $H_1$ are differentiable in the local regions considered.
In the passive mode, the represented target update has no direct action dependence:
\[
    \frac{\partial H_0(y)}{\partial a}=0.
\]
In the interaction mode, there exists a local interaction point
$(\bar u,\bar y,\bar a)$ and an action direction $v$ such that
\[
    D_a H_1(\bar u,\bar y,\bar a)v \neq 0.
\]
\end{assumption}

Assumption~\ref{ass:mode-dependent-controllability} implies that a target-only predictor is appropriate in the passive mode, but insufficient in the interaction mode.
The reason is that two anchors can share the same target state $y_t$ while differing in the actuated factor or action.
During interaction, these two anchors can produce different next target states.
A predictor depending only on $\psi_t$ would assign the same next representation to both anchors, and therefore cannot represent the interaction update.

\begin{lemma}[Need for an interaction-dependent residual]
\label{lem:need-affine-residual}
Suppose Assumption~\ref{ass:mode-dependent-controllability} holds.
In the passive mode, a first-order local predictor of the represented target update can be written as a target-only linear predictor.
In the interaction mode, no target-only linear predictor of the form $A\psi(y)$ can approximate the represented target update to first order in the interaction context.
Therefore, the interaction-mode local predictor requires an additional context-dependent residual, yielding the affine form
\[
    \psi_{t+1}
    =
    A_1\psi_t+b_t .
\]
\end{lemma}

\begin{proof}
We prove the two claims separately.

\textbf{Passive mode.}
In the passive mode,
\[
    \psi_{t+1}
    =
    H_0(y_t)
    =
    \psi(f_0(y_t)).
\]
Since $H_0$ depends only on $y_t$, its first-order local expansion in representation coordinates around $\bar y$ has the form
\[
    H_0(y)
    =
    H_0(\bar y)
    +
    A_0\bigl(\psi(y)-\psi(\bar y)\bigr)
    +
    o\!\left(\|\psi(y)-\psi(\bar y)\|\right).
\]
After centering the local coordinates, or equivalently absorbing the affine intercept into an augmented representation coordinate, this becomes
\[
    \psi_{t+1}
    =
    A_0\psi_t
    +
    o\!\left(\|\psi_t-\psi(\bar y)\|\right).
\]
Thus, a target-only linear predictor is sufficient to first order in the passive mode.

\textbf{Interaction mode.}
Assume for contradiction that a target-only linear predictor is first-order sufficient in the interaction mode.
Then there exists a matrix $A$ such that, for all sufficiently small $h$,
\[
    H_1(\bar u,\bar y,\bar a+hv)
    =
    A\psi(\bar y)
    +
    o(|h|).
\]
Evaluating the same expression at $h=0$ gives
\[
    H_1(\bar u,\bar y,\bar a)
    =
    A\psi(\bar y).
\]
Subtracting the two equations yields
\[
    H_1(\bar u,\bar y,\bar a+hv)
    -
    H_1(\bar u,\bar y,\bar a)
    =
    o(|h|).
\]
Dividing by $h$ and taking $h\rightarrow 0$ gives
\[
    D_aH_1(\bar u,\bar y,\bar a)v
    =
    0.
\]
This contradicts Assumption~\ref{ass:mode-dependent-controllability}, which states that
\[
    D_aH_1(\bar u,\bar y,\bar a)v \neq 0.
\]
Therefore, no target-only linear predictor can be first-order sufficient in the interaction mode.

It remains to show the affine form.
The first-order Taylor expansion of $H_1$ around $(\bar u,\bar y,\bar a)$ is
\[
    H_1(u,y,a)
    =
    H_1(\bar u,\bar y,\bar a)
    +
    A_1\bigl(\psi(y)-\psi(\bar y)\bigr)
    +
    D_uH_1(\bar u,\bar y,\bar a)(u-\bar u)
    +
    D_aH_1(\bar u,\bar y,\bar a)(a-\bar a)
    +
    o(\Delta),
\]
where
\[
    \Delta
    =
    \|\psi(y)-\psi(\bar y)\|
    +
    \|u-\bar u\|
    +
    \|a-\bar a\|.
\]
Rearranging the first-order terms gives
\[
    H_1(u,y,a)
    =
    A_1\psi(y)
    +
    b(u,y,a)
    +
    o(\Delta),
\]
where
\[
    b(u,y,a)
    =
    H_1(\bar u,\bar y,\bar a)
    -
    A_1\psi(\bar y)
    +
    D_uH_1(\bar u,\bar y,\bar a)(u-\bar u)
    +
    D_aH_1(\bar u,\bar y,\bar a)(a-\bar a).
\]
At time $t$, define
\[
    b_t \triangleq b(u_t,y_t,a_t).
\]
Therefore, the interaction-mode local predictor has the affine form
\[
    \psi_{t+1}
    =
    A_1\psi_t+b_t
    +
    o(\Delta).
\]
Dropping the higher-order local error gives the switched affine approximation used in the main analysis.
\end{proof}

% \begin{lemma}[Need for an interaction-dependent residual]
% \label{lem:need-affine-residual}
% Suppose Assumption~\ref{ass:mode-dependent-controllability} holds.
% In the passive mode, a local target predictor can be written as a target-only linear predictor.
% In the interaction mode, a target-only linear predictor is generally insufficient; an additional term depending on the interaction context is required.
% \end{lemma}

% \begin{proof}
% In the passive mode, $y_{t+1}=f_0(y_t)$.
% Therefore, after applying the target representation, the conditional next representation depends locally only on $\psi_t$.
% A local linear regression in representation space can therefore be written as
% \[
%     \psi_{t+1}\approx A_0\psi_t.
% \]

% In the interaction mode, $y_{t+1}=f_1(u_t,y_t,a_t)$.
% By Assumption~\ref{ass:mode-dependent-controllability}, changing $a_t$ while keeping $y_t$ fixed can change the next target state.
% Since $\psi$ is locally non-degenerate on goal-relevant target coordinates, this can also change $\psi_{t+1}$.
% Thus, two anchors with the same $\psi_t$ but different interaction contexts can require different predictions for $\psi_{t+1}$.
% No predictor of the form $A\psi_t$ can distinguish these two cases, because its input is identical.
% Hence, the interaction-mode predictor must include an additional interaction-dependent residual.
% \end{proof}

We now define the local affine approximation used in the proof.
The operator $A_0$ is the local linear predictor for passive target evolution.
The operator $A_1$ is the target-state coefficient in the interaction mode.
The residual $b_t$ is the part of the interaction update induced by the actuated factor and action after accounting for the target-state term.

\begin{assumption}[Mode-wise local affine representation dynamics]
\label{ass:mode-wise-linear}
Along the local trajectory segment analyzed in Appendix~\ref{app:one-reset-error}, the conditional mean dynamics of the target representation satisfy
\[
    \psi_{t+1}
    =
    A_0\psi_t,
    \qquad
    \omega_t=0,
\]
and
\[
    \psi_{t+1}
    =
    A_1\psi_t+b_t,
    \qquad
    \omega_t=1.
\]
Equivalently, the two cases can be written as the switched affine model
\[
    \psi_{t+1}
    =
    A_{\omega_t}\psi_t+\omega_t b_t.
\]
Here $b_t$ is zero outside interaction because it is multiplied by $\omega_t$.
\end{assumption}

This assumption is a local regression statement in representation space.
It does not require the original manipulation dynamics to be globally linear.
Instead, it says that near a trajectory segment, the passive target update is approximated by a target-only linear operator, while the interaction update requires a separate local operator plus an interaction-dependent residual.
This is the representation-level form of the dependency change in the factored dynamics.
The one-reset analysis below studies how an error in this residual is propagated to later endpoint energies.

\subsection{One-Reset Simplification}

We now formalize the one-reset setting used in Proposition~\ref{prop:endpoint-energy-error}.
Let $\tau$ denote the first future interaction offset after anchor time $t$.
The one-reset mode sequence is
\[
    \underbrace{0,\ldots,0}_{\tau\ \text{passive steps}},
    1,
    \underbrace{0,\ldots,0}_{k\ \text{passive steps}},
\]
where $k$ is the number of passive steps after the interaction.
Thus, the endpoint is at time $t+\tau+1+k$.

\begin{definition}[Interaction bridge shift]
\label{def:bridge-shift}
The interaction bridge shift is the difference between the actual interaction update and the passive
update that would have occurred without interaction:
\begin{equation}
    \delta_\tau
    \triangleq
    (A_1-A_0)\psi_{t+\tau}+b_{t+\tau}.
    \label{eq:bridge-shift-appendix}
\end{equation}
Equivalently,
\begin{equation}
    A_1\psi_{t+\tau}+b_{t+\tau}
    =
    A_0\psi_{t+\tau}
    +
    \delta_\tau .
    \label{eq:bridge-as-passive-plus-shift}
\end{equation}
\end{definition}

The bridge shift isolates the effect of the interaction itself.
Before the bridge, the target evolves passively.
At the bridge, the target receives the additional shift $\delta_\tau$.
After the bridge, this shift is carried forward by passive dynamics.

\begin{definition}[Bridge estimation error]
\label{def:bridge-error}
Let $\widehat{\delta}_\tau$ be the bridge shift represented by the learned energy model.
The local interaction-bridge estimation error is
\begin{equation}
    e
    \triangleq
    \widehat{\delta}_\tau-\delta_\tau .
    \label{eq:bridge-error-appendix}
\end{equation}
\end{definition}

\subsection{Future Mean Under One Interaction}

\begin{lemma}[One-reset future mean]
\label{lem:one-reset-future-mean}
Under Assumption~\ref{ass:mode-wise-linear} and the one-reset setting, the target representation $k$ passive steps after
the interaction has mean
\begin{equation}
    \mu_k
    =
    A_0^{\tau+1+k}\psi_t
    +
    A_0^k\delta_\tau .
    \label{eq:one-reset-future-mean}
\end{equation}
\end{lemma}

\begin{proof}
Before the interaction, the target follows passive dynamics for $\tau$ steps:
\begin{equation}
    \psi_{t+\tau}
    =
    A_0^\tau\psi_t .
    \label{eq:pre-interaction-passive}
\end{equation}
At the interaction step, Definition~\ref{def:bridge-shift} gives
\begin{equation}
    \psi_{t+\tau+1}
    =
    A_0\psi_{t+\tau}
    +
    \delta_\tau .
    \label{eq:interaction-update-with-shift}
\end{equation}
After the interaction, the target follows passive dynamics for $k$ more steps:
\begin{equation}
    \psi_{t+\tau+1+k}
    =
    A_0^k
    \left(
        A_0\psi_{t+\tau}
        +
        \delta_\tau
    \right).
    \label{eq:post-interaction-passive}
\end{equation}
Substituting Eq.~\eqref{eq:pre-interaction-passive} into Eq.~\eqref{eq:post-interaction-passive}
yields
\begin{align}
    \psi_{t+\tau+1+k}
    &=
    A_0^k
    \left(
        A_0A_0^\tau\psi_t
        +
        \delta_\tau
    \right) \\
    &=
    A_0^{\tau+1+k}\psi_t
    +
    A_0^k\delta_\tau .
\end{align}
This proves Eq.~\eqref{eq:one-reset-future-mean}.
\end{proof}

Lemma~\ref{lem:one-reset-future-mean} shows the key structure: the interaction bridge shift is
not observed at the endpoint as $\delta_\tau$, but as the propagated quantity $A_0^k\delta_\tau$.

\subsection{Propagation of Bridge Error}

\begin{lemma}[Propagated bridge error]
\label{lem:propagated-bridge-error}
If the learned bridge shift has error $e=\widehat{\delta}_\tau-\delta_\tau$, then the induced endpoint
mean error after $k$ passive steps is
\begin{equation}
    \widehat{\mu}_k-\mu_k
    =
    A_0^k e .
    \label{eq:future-mean-error}
\end{equation}
Consequently,
\begin{equation}
    \|\widehat{\mu}_k-\mu_k\|
    \le
    \|A_0^k\|\|e\|.
    \label{eq:future-mean-error-bound}
\end{equation}
\end{lemma}

\begin{proof}
By Lemma~\ref{lem:one-reset-future-mean}, the true endpoint mean is
\begin{equation}
    \mu_k
    =
    A_0^{\tau+1+k}\psi_t
    +
    A_0^k\delta_\tau .
    \label{eq:true-endpoint-mean}
\end{equation}
Using the learned bridge shift $\widehat{\delta}_\tau$ gives
\begin{equation}
    \widehat{\mu}_k
    =
    A_0^{\tau+1+k}\psi_t
    +
    A_0^k\widehat{\delta}_\tau .
    \label{eq:learned-endpoint-mean}
\end{equation}
Subtracting Eq.~\eqref{eq:true-endpoint-mean} from Eq.~\eqref{eq:learned-endpoint-mean} gives
\begin{align}
    \widehat{\mu}_k-\mu_k
    &=
    A_0^k
    \left(
        \widehat{\delta}_\tau-\delta_\tau
    \right) \\
    &=
    A_0^k e .
\end{align}
Taking norms and using submultiplicativity gives Eq.~\eqref{eq:future-mean-error-bound}.
\end{proof}

\subsection{Energy Error Induced by Bridge Error}

We now translate the propagated endpoint mean error into an energy error.
For the appendix proof, we use the whitened local Gaussian energy
\begin{equation}
    E_k
    =
    -\frac{1}{2}
    \left\|
        \psi(y_f)-\mu_k
    \right\|^2
    +
    \mathrm{const}.
    \label{eq:gaussian-energy-appendix}
\end{equation}
This is the isotropic form of the local Gaussian score used in the main analysis.
A non-isotropic covariance only changes the norm and introduces constant conditioning factors, which
are suppressed in the main-paper proportionality statement.

\begin{assumption}[Normalized local endpoint region]
\label{ass:normalized-local-region}
The sampled endpoint lies in the local quadratic region of the correct mean:
\begin{equation}
    \left\|
        \psi(y_f)-\mu_k
    \right\|
    \le
    1 .
    \label{eq:normalized-local-region}
\end{equation}
\end{assumption}

This normalization is only for presentation.
If the local residual is bounded by a constant other than $1$, the first term in the final bound is
multiplied by that constant.

\begin{lemma}[Exact energy perturbation]
\label{lem:exact-energy-perturbation}
Suppose the propagated bridge-error relation in Lemma~\ref{lem:propagated-bridge-error} holds.
Then the energy perturbation caused by bridge error is
\begin{equation}
    \widehat{E}_k-E_k
    =
    \left(
        \psi(y_f)-\mu_k
    \right)^\top
    A_0^k e
    -
    \frac{1}{2}
    \left\|
        A_0^k e
    \right\|^2 .
    \label{eq:energy-perturbation-appendix}
\end{equation}
\end{lemma}

\begin{proof}
By Lemma~\ref{lem:propagated-bridge-error},
\begin{equation}
    \widehat{\mu}_k
    =
    \mu_k+A_0^k e .
    \label{eq:mean-with-propagated-error}
\end{equation}
Using Eq.~\eqref{eq:gaussian-energy-appendix}, the learned energy is
\begin{equation}
    \widehat{E}_k
    =
    -\frac{1}{2}
    \left\|
        \psi(y_f)-\widehat{\mu}_k
    \right\|^2
    +
    \mathrm{const}.
\end{equation}
Substituting Eq.~\eqref{eq:mean-with-propagated-error} gives
\begin{equation}
    \widehat{E}_k
    =
    -\frac{1}{2}
    \left\|
        \psi(y_f)-\mu_k-A_0^k e
    \right\|^2
    +
    \mathrm{const}.
\end{equation}
Therefore,
\begin{align}
    \widehat{E}_k-E_k
    &=
    -\frac{1}{2}
    \left[
        \left\|
            \psi(y_f)-\mu_k-A_0^k e
        \right\|^2
        -
        \left\|
            \psi(y_f)-\mu_k
        \right\|^2
    \right] \\
    &=
    \left(
        \psi(y_f)-\mu_k
    \right)^\top
    A_0^k e
    -
    \frac{1}{2}
    \left\|
        A_0^k e
    \right\|^2 .
\end{align}
This proves Eq.~\eqref{eq:energy-perturbation-appendix}.
\end{proof}

\begin{corollary}[Endpoint energy-error bound]
\label{cor:endpoint-energy-error-bound}
Suppose Lemma~\ref{lem:exact-energy-perturbation} holds and Assumption~\ref{ass:normalized-local-region} holds.
Then
\begin{equation}
    |\widehat{E}_k-E_k|
    \le
    \|A_0^k\|\|e\|
    +
    \frac{1}{2}
    \|A_0^k\|^2\|e\|^2 .
    \label{eq:energy-error-bound-appendix}
\end{equation}
\end{corollary}

\begin{proof}
Starting from Lemma~\ref{lem:exact-energy-perturbation}, apply the triangle inequality:
\begin{align}
    |\widehat{E}_k-E_k|
    &\le
    \left|
        \left(
            \psi(y_f)-\mu_k
        \right)^\top
        A_0^k e
    \right|
    +
    \frac{1}{2}
    \left\|
        A_0^k e
    \right\|^2 .
    \label{eq:triangle-bound}
\end{align}
By Cauchy--Schwarz,
\begin{equation}
    \left|
        \left(
            \psi(y_f)-\mu_k
        \right)^\top
        A_0^k e
    \right|
    \le
    \left\|
        \psi(y_f)-\mu_k
    \right\|
    \left\|
        A_0^k e
    \right\|.
\end{equation}
Using Assumption~\ref{ass:normalized-local-region},
\begin{equation}
    \left|
        \left(
            \psi(y_f)-\mu_k
        \right)^\top
        A_0^k e
    \right|
    \le
    \left\|
        A_0^k e
    \right\|.
\end{equation}
Finally, by submultiplicativity,
\begin{equation}
    \left\|
        A_0^k e
    \right\|
    \le
    \|A_0^k\|\|e\|.
\end{equation}
Substituting these bounds into Eq.~\eqref{eq:triangle-bound} gives
\begin{equation}
    |\widehat{E}_k-E_k|
    \le
    \|A_0^k\|\|e\|
    +
    \frac{1}{2}
    \|A_0^k\|^2\|e\|^2 .
\end{equation}
\end{proof}

Corollary~\ref{cor:endpoint-energy-error-bound} is the formal version of Proposition~\ref{prop:endpoint-energy-error}.
Removing constants associated with local normalization and whitening gives the main-paper scaling:
\[
    \sup |\widehat{E}_k-E_k|
    \propto
    \|A_0^k\|\|e\|
    +
    \frac{1}{2}
    \|A_0^k\|^2\|e\|^2 .
\]

\subsection{Interpretation}

The proof separates the endpoint error into two factors.
The first factor, $\|e\|$, is the local error in estimating the interaction bridge shift.
This error becomes large when few training samples directly cover the interaction boundary.
The second factor, $\|A_0^k\|$, is the passive post-interaction propagation factor.
Even if the bridge error is small, passive dynamics can rotate, attenuate, or amplify the error before
the endpoint is used in the contrastive objective.

This explains why endpoint-only CRL can learn fragile energies near contact.
If the positive future is sampled long after the interaction, CRL supervises the interaction action only
through the propagated term $A_0^k\delta_\tau$.
Interaction-weighted resampling targets the two terms in the bound: it increases bridge-centered
samples, reducing $\|e\|$, and it favors endpoints closer to the interaction, reducing the effective
post-interaction propagation length $k$.